\section{Precedents in Grothendieck's notes}\label{sec:precedents}

This section collects two precedents in Grothendieck's notes. The decks of his notes are in the registry (\S\ref{sec:grdecks}).

\subsection{Preorders and spatial topoi: folder 161-3, p.~26, and folder 37, p.~48}

In an undated manuscript note on topoi \cite{grothendieck-161-3}, p.~26, Grothendieck lists equivalent conditions on a topos $E$: to be equivalent both to the topos of a space and to a presheaf topos; to be the topos of an \q{essential} space (a space all of whose points are essential); to be equivalent to $\widehat I$ for an ordered set $I$; to be generated by the subobjects of the final object and to have enough essential points. He asks for a proof that $\widehat C$ satisfies these conditions iff $C$ is preordered, and that the topos of a space $X$ does iff the essential points of $X$ are very dense; he also notes 2-equivalences between essential spatial topoi, preorders and essential spaces. In our terms these are bridges on the deck of preorders, with the arches of subterminal objects and of essential points. The generator re-derives the first exercise from the Morita arch: if $\widehat C\simeq\widehat P$ for a preorder $P$, then $\Kar(C)\simeq\Kar(P)\simeq P$, and $C$, a full subcategory of its Karoubi envelope, is thin; conversely a thin category is equivalent to its thin skeleton, a partial order (Lean: \lean{exists_preorder_moritaEquivalentCat_iff_isThin}). The second exercise, on spaces, is not formalised here. The same identification appears in his notes for the re-edition of SGA \cite[p.~48]{grothendieck-37}: for an ordered set $I$ the opens of $\widehat I$ are the sieves of $I$, hence $\widehat I\simeq\mathrm{Top}(I)$.

\subsection{Subtopoi and quotients of a theory: folder 161-2, pp.~17 and 44}

In his notes on universal algebra \cite[p.~17]{grothendieck-161-2}, under the heading \q{sous-théorie d'une $\Delta$-théorie (représentable par un topos) et sous-topos}, Grothendieck takes a topos $R$ and a set $\Sigma$ of arrows of $R$: the morphisms that invert the arrows of $\Sigma$ factor through a subtopos $R'$ of $R$, and every subtopos of $R$ arises in this way. Read with $R$ the classifying topos of a theory, this is in substance the correspondence between quotients of a theory and subtopoi of its classifying topos on which Caramello builds her bridges \cite{caramello-tst}. On p.~44 of the same folder he reads structures as open subtopoi of a classifying topos.
